../cictikz/src/cictikz/data/tex/cictikz_preamble.tex

% The tree decode, drawn as pass devices rather than abstract circles:
% every tap gets a switch steered by b_0 or its complement, every joined
% pair needs another switch behind it, and the doubling repeats until one
% node is left. Drawn for a three-bit string: 8 + 4 + 2 = 14 switches is
% the 2^(N+1) - 2 the lecture counts, and each rank's control label shows
% which bit pays for it. Same NMOS-with-control language as dac_r_div*
% and dac_r_rows, so the complexity slides read as one family.

\begin{circuitikz}[american, thick, transform shape, circuit ee IEC]

  ../cictikz/src/cictikz/data/tex/cictikz_lib.tex

  \def\ry{1.9}       % leaf pitch

  % ---- Rank 0: one switch per tap, b_0 on odd taps ----
  \foreach \k in {0,...,7} {
    \draw (0,{\k*\ry}) \portIn{$V_\k$}
          -- ++(0.36,0) to [Tnmos, n=A\k] ++(1.0,0)
          -- ++(0.3,0) coordinate (oA\k);
    \draw (A\k.gate) -- ++(0,0.35) coordinate (gA\k);
    \pgfmathparse{int(mod(\k,2))}
    \ifnum\pgfmathresult=1
      \node[anchor=south] at (gA\k) {$b_0$};
    \else
      \node[anchor=south] at (gA\k) {$\overline{b_0}$};
    \fi
  }

  % ---- Rank 1: each pair joins into one switch, b_1 decides ----
  \foreach \m in {0,...,3} {
    \pgfmathsetmacro{\ym}{(2*\m+0.5)*\ry}
    \coordinate (eB\m) at (3.2,\ym);
    \pgfmathtruncatemacro{\ka}{2*\m}
    \pgfmathtruncatemacro{\kb}{2*\m+1}
    \draw (oA\ka) -- (eB\m);
    \draw (oA\kb) -- (eB\m);
    \fill (eB\m) circle (0.07);
    \draw (eB\m) to [Tnmos, n=B\m] ++(1.0,0) -- ++(0.3,0) coordinate (oB\m);
    \draw (B\m.gate) -- ++(0,0.35) coordinate (gB\m);
    \pgfmathparse{int(mod(\m,2))}
    \ifnum\pgfmathresult=1
      \node[anchor=south] at (gB\m) {$b_1$};
    \else
      \node[anchor=south] at (gB\m) {$\overline{b_1}$};
    \fi
  }

  % ---- Rank 2, and the output the last pair joins into ----
  \foreach \m in {0,1} {
    \pgfmathsetmacro{\ym}{(4*\m+1.5)*\ry}
    \coordinate (eC\m) at (6.2,\ym);
    \pgfmathtruncatemacro{\ka}{2*\m}
    \pgfmathtruncatemacro{\kb}{2*\m+1}
    \draw (oB\ka) -- (eC\m);
    \draw (oB\kb) -- (eC\m);
    \fill (eC\m) circle (0.07);
    \draw (eC\m) to [Tnmos, n=C\m] ++(1.0,0) -- ++(0.3,0) coordinate (oC\m);
    \draw (C\m.gate) -- ++(0,0.35) coordinate (gC\m);
    \ifnum\m=1
      \node[anchor=south] at (gC\m) {$b_2$};
    \else
      \node[anchor=south] at (gC\m) {$\overline{b_2}$};
    \fi
  }

  \coordinate (vo) at (9.2,{3.5*\ry});
  \draw (oC0) -- (vo);
  \draw (oC1) -- (vo);
  \fill (vo) circle (0.07);
  \draw (vo) -- ++(0.5,0) coordinate (vp);
  \draw (vp) circle (0.06);
  \node[anchor=west, xshift=3pt] at (vp) {$V_{OUT}$};

\end{circuitikz}
\end{document}
